\Needspace{4\baselineskip}
\begin{Example}
\label{ex:sobolev-first-second-norms}
对 $1\leq p<\infty$, 有\[
\|u\|_{1,p}=\left\{\int_{\Omega}\left(|u|^p+\sum_{i=1}^m\left|\dfrac{\partial u}{\partial x_i}\right|^p\right)\,\mathrm{d}\bm{x}\right\}^{1/p},
\]
\[
\|u\|_{2,p}=\left\{\int_{\Omega}\left(|u|^p+\sum_{i=1}^m\left|\dfrac{\partial u}{\partial x_i}\right|^p+\sum_{i,j=1}^m\left|\dfrac{\partial^2u}{\partial x_i\partial x_j}\right|^p\right)\,\mathrm{d}\bm{x}\right\}^{1/p}.
\]
\end{Example}

\begin{Example}
\label{ex:heaviside-absolute-value-sobolev-membership}
Heaviside 函数 $H(x)$ 对任意 $p\in[1,\infty]$ 都属于 $L^p(-1,1)$. 然而, 它不属于 $W^{1,p}(-1,1)$, 因为其分布导数, 即 Dirac delta, 不能由 $L^p(-1,1)$ 函数表示. 类似地, 绝对值函数 $f(x)=|x|$ 属于 $W^{1,p}(-1,1)$, 因为其分布导数可以由以下 $L^p(-1,1)$ 函数表示:\[f'(x)=\begin{cases}-1&x<0\\1&x>0\end{cases}\]然而, $f$ 不属于 $W^{2,p}(-1,1)$, 因为其二阶分布导数 $(2\delta)$ 不能由 $L^p$ 中的函数表示.
\end{Example}

\begin{Example}
\label{ex:radial-power-sobolev-membership}
考察函数\[f(\bm{x}):=|\bm{x}|^{\alpha}=\left(\sum_{i=1}^m x_i^2\right)^{\alpha/2}\]定义在 $B=B_1(\bm{0})\subset\mathbb{R}^m$ 上, 这是 $\mathbb{R}^m$ 中半径为一的球. 利用球坐标很容易检验, 若 $\alpha+m>0$, 则此函数可积. 因而若 $p\alpha>-m$, 则有 $f\in L^p(B)$.

形式地取偏导数得到\[\dfrac{\partial f}{\partial x_i}(\bm{x})=\alpha x_i\left(\sum_{i=1}^m x_i^2\right)^{\alpha/2-1}=\alpha|\bm{x}|^{(\alpha-2)/2}x_i\]\[
\dfrac{\partial^2f}{\partial x_i\partial x_j}(\bm{x})=\alpha(\alpha-2)x_ix_j\left(\sum_{i=1}^m x_i^2\right)^{\alpha/2-2}+\alpha\left(\sum_{i=1}^m x_i^2\right)^{\alpha/2-1}\delta_{ij}
\]其中 $\delta_{ij}$ 为 Kronecker delta.

可用反常积分的标准方法证明, 若形式导数具有可积奇点, 则该函数为分布导数的一个代表元. 利用这一点并计算范数, 经过一些工作得到\begin{itemize}
\item $f\in W^{1,p}(B)$ 当且仅当 $p(\alpha-1)>-m$ (即 $\alpha>1-m/p$).
\item $f\in W^{2,p}(B)$ 当且仅当 $p(\alpha-2)>-m$ (即 $\alpha>2-m/p$).
\end{itemize}类似的计算表明, $|\bm{x}|^{\alpha}\in W^{k,p}(B)$ 当且仅当 $\alpha>k-m/p$.
\end{Example}
